% !TeX program = pdflatex
\documentclass[11pt,a4paper]{amsart}

\newcommand{\notelanguage}{finnish}
% Shared preamble for course notes and exercise sets.

\usepackage[T1]{fontenc}
\usepackage{lmodern}
\providecommand{\notelanguage}{english}
\usepackage[finnish,english]{babel}
\AtBeginDocument{\expandafter\selectlanguage\expandafter{\notelanguage}}

\newcommand{\notesfinnishlanguage}{finnish}
\ifx\notelanguage\notesfinnishlanguage
\newcommand{\notestheoremname}{Lause}
\newcommand{\noteslemmaname}{Lemma}
\newcommand{\notespropositionname}{Propositio}
\newcommand{\notescorollaryname}{Korollaari}
\newcommand{\notesdefinitionname}{Määritelmä}
\newcommand{\notesexamplename}{Esimerkki}
\newcommand{\notesproblemname}{Tehtävä}
\newcommand{\notesexercisename}{Tehtävä}
\newcommand{\notesremarkname}{Huomautus}
\newcommand{\notessolutionname}{Ratkaisu}
\else
\newcommand{\notestheoremname}{Theorem}
\newcommand{\noteslemmaname}{Lemma}
\newcommand{\notespropositionname}{Proposition}
\newcommand{\notescorollaryname}{Corollary}
\newcommand{\notesdefinitionname}{Definition}
\newcommand{\notesexamplename}{Example}
\newcommand{\notesproblemname}{Problem}
\newcommand{\notesexercisename}{Exercise}
\newcommand{\notesremarkname}{Remark}
\newcommand{\notessolutionname}{Solution}
\fi

\usepackage[a4paper,margin=30mm]{geometry}
\usepackage{microtype}
\usepackage{graphicx}
\usepackage{enumitem}

\newcommand{\exerciseheading}[2]{%
  \par\noindent
  {\large\bfseries #1\par}
  \vspace{0.1em}
  \noindent{\large #2\par}
  \vspace{1.5em}
}

\usepackage{amsmath}
\usepackage{amssymb}
\usepackage{amsthm}
\usepackage{mathtools}
\usepackage{aliascnt}

% Row operations beside the affected row, e.g. \rowop{+2R_1}.
\newcommand{\rowop}[1]{\mathrel{\big|}\,#1}

% Linear systems: the argument gives the number of alignment pairs.
\newenvironment{linearsystem}[1]
{\left\{
  \begin{alignedat}{#1}}
    {
  \end{alignedat}\right.}

  \usepackage[hidelinks,pdfusetitle]{hyperref}

  \numberwithin{equation}{section}
  \setlist{itemsep=0.25em,topsep=0.5em}

  \theoremstyle{plain}
  \newtheorem{theorem}{\notestheoremname}[section]
  \newaliascnt{lemma}{theorem}
  \newtheorem{lemma}[lemma]{\noteslemmaname}
  \aliascntresetthe{lemma}
  \newaliascnt{proposition}{theorem}
  \newtheorem{proposition}[proposition]{\notespropositionname}
  \aliascntresetthe{proposition}
  \newaliascnt{corollary}{theorem}
  \newtheorem{corollary}[corollary]{\notescorollaryname}
  \aliascntresetthe{corollary}

  \theoremstyle{definition}
  \newaliascnt{definition}{theorem}
  \newtheorem{definition}[definition]{\notesdefinitionname}
  \aliascntresetthe{definition}
  \newaliascnt{example}{theorem}
  \newtheorem{example}[example]{\notesexamplename}
  \aliascntresetthe{example}
  \newaliascnt{problem}{theorem}
  \newtheorem{problem}[problem]{\notesproblemname}
  \aliascntresetthe{problem}
  \newtheorem{exercise}{\notesexercisename}

  \theoremstyle{remark}
  \newaliascnt{remark}{theorem}
  \newtheorem{remark}[remark]{\notesremarkname}
  \aliascntresetthe{remark}

  \newenvironment{solution}
  {
  \begin{proof}[\notessolutionname]}
    {
  \end{proof}}

  \usepackage[nameinlink,noabbrev]{cleveref}

  \crefname{theorem}{theorem}{theorems}
  \crefname{lemma}{lemma}{lemmas}
  \crefname{proposition}{proposition}{propositions}
  \crefname{corollary}{corollary}{corollaries}
  \crefname{definition}{definition}{definitions}
  \crefname{example}{example}{examples}
  \crefname{problem}{problem}{problems}
  \crefname{exercise}{exercise}{exercises}
  \crefname{remark}{remark}{remarks}

  \ifx\notelanguage\notesfinnishlanguage
  \crefname{theorem}{lause}{lauseet}
  \crefname{lemma}{lemma}{lemmat}
  \crefname{proposition}{propositio}{propositiot}
  \crefname{corollary}{korollaari}{korollaarit}
  \crefname{definition}{määritelmä}{määritelmät}
  \crefname{example}{esimerkki}{esimerkit}
  \crefname{problem}{tehtävä}{tehtävät}
  \crefname{exercise}{tehtävä}{tehtävät}
  \crefname{remark}{huomautus}{huomautukset}
  \fi

  \newcommand{\NN}{\mathbb{N}}
  \newcommand{\NNone}{\mathbb{N}_1}
  \newcommand{\ZZ}{\mathbb{Z}}
  \newcommand{\QQ}{\mathbb{Q}}
  \newcommand{\RR}{\mathbb{R}}
  \newcommand{\CC}{\mathbb{C}}
  \newcommand{\PP}{\mathbb{P}}

  \DeclareMathOperator{\im}{im}
  \DeclareMathOperator{\rank}{rank}
  \DeclareMathOperator{\syt}{syt}
  \DeclareMathOperator{\pyj}{pyj}
  \DeclareMathOperator{\dom}{dom}
  \DeclarePairedDelimiter{\abs}{\lvert}{\rvert}
  \DeclarePairedDelimiter{\norm}{\lVert}{\rVert}
  \newcommand{\set}[1]{\left\lbrace#1\right\rbrace}
  \newcommand{\maxset}[1]{\max\set{#1}}
  \newcommand{\minset}[1]{\min\set{#1}}
  \newcommand{\powerset}[1]{\mathcal{P}\!\left(#1\right)}


\title{Analyysi II\\Harjoitus 2}
\author{}
\date{}

\begin{document}

\exerciseheading{Analyysi II}{Harjoitus 2}

\begin{exercise}
  Bolzanon lauseen 4.2.1 todistuksessa todistettiin vain tapaus
  \(f(a) < 0\) ja \(f(b) > 0\). Todista tapaus \(f(a) > 0\) ja \(f(b) < 0\)
  tämän avulla ja tutkimalla funktiota \(-f\).
\end{exercise}

\begin{solution}
  Olkoon \(g \colon [a, b] \to \RR\), \(g(x) = -f(x)\). Koska funktiot \(f\) ja
  \(x \mapsto -1\) ovat jatkuvia, on myös niiden tulo \(g\) jatkuva.
  Oletetaan, että \(f(a) > 0\) ja \(f(b) < 0\). Tällöin \(g(a) < 0\) ja
  \(g(b) > 0\). Funktio \(g\) täten täyttää jo todistetun tapauksen ehdot,
  eli on olemassa \(x \in (a, b)\) siten, että \(g(x) = 0\). Lisäksi
  \[
    f(x) = -g(x) = 0.
  \]
  Täten tapaus \(f(a) > 0\) ja \(f(b) < 0\) on todistettu.

\end{solution}

\begin{exercise}
  Osoita, että yhtälöllä \(x^7 = x + 7\) on ainakin yksi ratkaisu reaalilukujen
  joukossa.
\end{exercise}

\begin{solution}
  Olkoon \(f \colon \RR \to \RR\), \(f(x) = x^7 - x - 7\). Huomataan, että
  \(f(0) = -7 < 0\) ja \(f(2) = 119 > 0\). Funktio \(f\) on polynomifunktiona
  jatkuva, joten Bolzanon lauseen nojalla on olemassa \(x \in (0, 2)\)
  siten, että \(f(x) = 0\). Tämä \(x\) on myös alkuperäisen yhtälön
  \(x^7 = x + 7\) ratkaisu.
\end{solution}

\begin{exercise}
  Rajoita funktio \(f(x) = (x^3 - 1)^2\) väleihin, joilla se on aidosti
  monotoninen, ja määritä vastaavat käänteisfunktiot.
\end{exercise}

\begin{solution}
  Olkoon \(g \colon \RR \to \RR\), \(g(x) = x^3 - 1\). Kun
  \(x, y \in \RR\) ja \(x < y\), saadaan
  \[
    \begin{aligned}
      g(y) - g(x) &= y^3 - x^3 \\
      &= (y - x)(y^2 + yx + x^2).
    \end{aligned}
  \]
  Tässä \(y - x > 0\) ja
  \[
    y^2 + yx + x^2 = \left(y + \frac{x}{2}\right)^2 + \frac{3}{4}x^2 > 0.
  \]
  Jälkimmäinen lauseke voisi olla nolla vain, jos \(x = y = 0\), mikä on
  ristiriidassa ehdon \(x < y\) kanssa. Siis \(g(x) < g(y)\), joten \(g\)
  on aidosti kasvava.
  Koska \(g(1) = 0\), on \(g(x) \leq 0\), kun \(x \leq 1\), ja
  \(g(x) \geq 0\), kun \(x \geq 1\).
  Neliöfunktio on aidosti vähenevä välillä \((-\infty, 0]\) ja aidosti
  kasvava välillä \([0, \infty)\). Näin ollen \(f(x) = [g(x)]^2\) on
  aidosti vähenevä välillä \((-\infty, 1]\) ja aidosti kasvava välillä
  \([1, \infty)\).

  Merkitään näitä rajoittumia \(f_- = f|_{(-\infty, 1]}\) ja
  \(f_+ = f|_{[1, \infty)}\). Määritetään niiden käänteisfunktiot.
  Olkoon \(y \geq 0\). Kun \(x \leq 1\), saadaan
  \[
    \begin{aligned}
      y = (x^3 - 1)^2
      &\iff \sqrt{y} = 1 - x^3 \\
      &\iff x = \sqrt[3]{1 - \sqrt{y}}.
    \end{aligned}
  \]
  Vastaavasti, kun \(x \geq 1\), saadaan
  \[
    \begin{aligned}
      y = (x^3 - 1)^2
      &\iff \sqrt{y} = x^3 - 1 \\
      &\iff x = \sqrt[3]{1 + \sqrt{y}}.
    \end{aligned}
  \]
  Jokaisella \(y \geq 0\) ensimmäinen kaava antaa arvon \(x \leq 1\) ja
  toinen arvon \(x \geq 1\). Kummankin rajoittuman arvojoukko on siis
  \([0, \infty)\), ja aidon monotonisuuden vuoksi ne ovat injektioita.
  Täten vastaavat käänteisfunktiot ovat
  \[
    \begin{aligned}
      f_-^{-1} &\colon [0, \infty) \to (-\infty, 1],
      & f_-^{-1}(y) &= \sqrt[3]{1 - \sqrt{y}}, \\
      f_+^{-1} &\colon [0, \infty) \to [1, \infty),
      & f_+^{-1}(y) &= \sqrt[3]{1 + \sqrt{y}}.
    \end{aligned}
    \qedhere
  \]
\end{solution}

\begin{exercise}
  Anna esimerkki funktioista \(f, g \colon (0, 1) \to (0, 1)\), jotka eivät
  ole jatkuvia, ja joille \(g \circ f \colon (0, 1) \to (0, 1)\) on jatkuva
  ja \(f \circ g \colon (0, 1) \to (0, 1)\) ei ole jatkuva. Esimerkkisi
  osoittaa, että kahden epäjatkuvan funktion yhdistetty funktio voi olla
  jatkuva tai epäjatkuva.
\end{exercise}

\begin{solution}
  Määritellään funktiot \(f, g \colon (0, 1) \to (0, 1)\) kaavoilla
  \[
    f(x) =
    \begin{cases}
      \frac{3}{4}, & x = \frac{1}{2}, \\
      \frac{1}{4}, & x \in (0, 1) \setminus \{\frac{1}{2}\},
    \end{cases}
    \qquad \text{ja} \qquad
    g(x) =
    \begin{cases}
      \frac{1}{2}, & x = \frac{1}{2}, \\
      \frac{1}{4}, & x \in (0, 1) \setminus \{\frac{1}{2}\}.
    \end{cases}
  \]
  Molempien funktioiden arvot kuuluvat välille \((0, 1)\).
  Kun \(x \ne \frac{1}{2}\), on \(f(x) = g(x) = \frac{1}{4}\), joten
  \[
    \lim_{x \to 1/2} f(x) = \lim_{x \to 1/2} g(x) = \frac{1}{4}.
  \]
  Koska \(f\left(\frac{1}{2}\right) = \frac{3}{4}\) ja
  \(g\left(\frac{1}{2}\right) = \frac{1}{2}\), kumpikaan funktio ei ole jatkuva
  pisteessä \(\frac{1}{2}\).

  Funktio \(f\) saa vain arvot \(\frac{1}{4}\) ja \(\frac{3}{4}\), joille
  \(g\left(\frac{1}{4}\right) = g\left(\frac{3}{4}\right) =
  \frac{1}{4}\). Näin ollen
  \[
    (g \circ f)(x) = \frac{1}{4}
    \qquad \text{kaikilla } x \in (0, 1),
  \]
  joten \(g \circ f\) on vakiofunktiona jatkuva.

  Toisaalta \(f(g\left(\frac{1}{2}\right)) = f\left(\frac{1}{2}\right) =
  \frac{3}{4}\), ja kun
  \(x \ne \frac{1}{2}\), saadaan
  \(f(g(x)) = f\left(\frac{1}{4}\right) = \frac{1}{4}\). Siis
  \[
    (f \circ g)(x) =
    \begin{cases}
      \frac{3}{4}, & x = \frac{1}{2}, \\
      \frac{1}{4}, & x \in (0, 1) \setminus \{\frac{1}{2}\},
    \end{cases}
    = f(x).
  \]
  Täten \(f \circ g\) ei ole jatkuva pisteessä \(\frac{1}{2}\).
\end{solution}

\begin{exercise}
  Olkoon \(n \in \NNone\) pariton. Osoita, että funktiolla
  \(f \colon \RR \to \RR\), \(f(x) = x^n\), on jatkuva ja aidosti kasvava
  käänteisfunktio \(f^{-1} \colon \RR \to \RR\).
\end{exercise}

\begin{solution}
  Olkoon \(n \in \NNone\) pariton. Jos \(n = 1\), niin \(f(x) = x\),
  joten väite on selvä. Oletetaan siis, että \(n \geq 3\).
  Olkoon \(x, y \in \RR\) siten, että
  \(x < y\). Kun \(0 \leq x < y\), niin
  \[
    y^n - x^n = (y - x)(y^{n - 1} + y^{n - 2}x + \cdots
    + x^{n - 1}) > 0,
  \]
  sillä \(y - x > 0\), summan kaikki termit ovat epänegatiivisia ja
  \(y^{n - 1} > 0\). Siis \(f(x) < f(y)\).
  Kun \(x < y \leq 0\), niin \(0 \leq (-y) < (-x)\).
  Edellisen tapauksen ja luvun \(n\) parittomuuden nojalla
  \[
    (-y)^n < (-x)^n \iff x^n < y^n,
  \]
  eli \(f(x) < f(y)\). Lopuksi, kun \(x < 0 < y\), niin luvun \(n\)
  parittomuuden nojalla
  \[
    x^n < 0 < y^n.
  \]
  On osoitettu, että \(f\) on aidosti kasvava. Lisäksi \(f\) on jatkuva
  ollessaan polynomifunktio.

  Olkoon \(a \in \RR\) ja \(M = \abs{a} + 1\). Koska \(n \geq 1\), niin
  \(M^n \geq M\), joten \(n\) parittomuuden nojalla
  \[
    f(-M) = -M^n \leq -M < a < M \leq M^n = f(M).
  \]
  Nyt väliarvolauseen (korollaari 4.2.2) nojalla on olemassa
  \(x \in (-M, M)\) siten, että \(f(x) = a\). Koska \(a\) oli
  mielivaltainen reaaliluku, saadaan \(f(\RR) = \RR\).
  Nyt lauseen 4.2.4 nojalla funktiolla \(f\) on jatkuva ja aidosti kasvava
  käänteisfunktio \(f^{-1} \colon \RR \to \RR\).
\end{solution}

\end{document}
